\section{Scope and relationship to the original study}
\label{sec:scope}

This study is an adaptation of the ESG-washing index of \citet{lagasio2024}, not a
strict replication of it. We re-implement the index architecture and examine its
measurement properties, but three inputs differ by design --- the corpus, the ESG
vocabulary and the sentiment dictionary --- and one of the three cannot be aligned at
all from the published record, because the original never publishes its vocabulary or
its size. The delta is stated
here, at the outset, so that no later result is read as a failed reproduction of
something it never attempted to reproduce.

\paragraph{What is carried over.} Two things. First, the index architecture: a
document-level tone score and a document-level substance score, each rescaled across
the sample, differenced, and thresholded at zero to produce a binary label. Second,
the tone-minus-substance logic that motivates it --- positive language in excess of
measurable substantive content as the observable signature of ESG-washing. Both are
carried over unchanged, including the zero threshold and the labelling rule, which the
original states explicitly (\S\ref{sec:index}). Where the original's specification is incomplete
we do not close the gap silently: \S\ref{sec:original} documents each silence, and
\S\ref{sec:index} carries three substance specifications precisely because the
original's description does not determine one.

\paragraph{Where we diverge, and why.} Table~\ref{tab:scope-divergence} sets out the
full delta. Three of its entries are input choices that require a reason; a fourth is a
pipeline stage with no counterpart in the original.

\emph{Corpus.} The original starts from the top 2{,}000 firms by market capitalisation
in Refinitiv and retains the 749 with a standalone, English-language ESG report for
2023; ``companies that integrated their ESG disclosure within their annual financial
reports were excluded'' \citep[p.~3]{lagasio2024}. Our corpus is drawn largely from the population that
rule excludes: 343 unique reports from 49 firms in seven European countries over
2018--2024, of which 235 are annual reports, 101 universal registration documents (URD)
and 7 are standalone sustainability reports. The entailment runs from the panel design,
and we state its direction rather than present the divergence as unavoidable: we adopt a
balanced seven-year panel --- longitudinal work being what the original's own conclusion
calls for --- and it is that choice which rules out the original's inclusion criterion. A
balanced panel requires one annually recurring disclosure per firm in every year of the
window. The annual report, and in France the URD, is a periodic regulatory filing that
every firm in the frame publishes in every year; a standalone sustainability report is
issued at the firm's discretion, so neither its existence nor its recurrence over seven
consecutive years can be assumed. Our own corpus shows how binding that is here: of the 49
firms, exactly one --- Sanofi-Aventis --- is represented by standalone sustainability
reports across the window, and it is so in all seven years, so a panel restricted to that
document type would have contained one firm. What our records hold is the document
\emph{selected} for each firm-year, not what each firm made available, so this is a
statement about what a document-type restriction would have left us with and not about what
these firms published. Relatedly, 43 of the 49 firms are represented by a single document
type across the whole window; the six that change do so once, from annual report to URD,
consistent with the French regime. Retaining the original's rule would have made panel
membership a function of the firm's chosen reporting vehicle --- a disclosure behaviour
in its own right --- and left the panel unbalanced. The restriction to 49 firms in seven
European countries is a separate sampling-frame decision and not a consequence of the
panel design; it, and the operational consequence of mixing document types inside a
single index, are treated in \S\ref{sec:data} and \S\ref{sec:limitations}.

\emph{Lexicon.} The original's keyword list is attributed to GRI (2021) and SASB
(2018), but its size is never stated, the list is never published, and the data
availability statement reads in full ``No'' (\S\ref{sec:original}). We therefore built
a purpose-designed 154-term vocabulary and publish it in full (\S\ref{sec:repro}). The
consequence is stated plainly rather than deferred to a limitations paragraph: our
substance score and the original's are not the same measurement, and no numerical
difference between the two indices can be attributed cleanly to any single design
choice (\S\ref{sec:lexicon}).

\emph{Sentiment dictionary.} The original scores tone with TextBlob polarity, a
general-purpose pattern-based tool. We use the Loughran--McDonald finance dictionary
\citep{loughran2011}, which is calibrated on disclosure language; the original's own
conclusion names ``refining sentiment analysis techniques'' \citep[p.~10]{lagasio2024} among the extensions
it anticipates. The two are not interchangeable instruments, and each carries its own
defects --- ours are set out in \S\ref{sec:limitations}, and what the original discloses
about its own instrument is audited in \S\ref{sec:original} --- so the tone components
are no more directly comparable across the two studies than the substance components
are.

\emph{Text unit.} A fourth divergence is a stage rather than an input: we vectorise
lexically extracted ESG zones, not whole reports (\S\ref{sec:extraction}). It has no
counterpart in the original, which vectorises the report entire; it is required here
because ESG disclosure in an annual report or URD sits inside a much larger financial
and legal document; and it fixes the denominator against which the length effects of
\S\ref{sec:measurement} operate. Its cost --- an untested filtering stage upstream of
every number we report --- is conceded in \S\ref{sec:limitations}.

\paragraph{Temporal extension.} The original states
that a panel was unnecessary for its purpose --- ``a cross-sectional dataset focused on
a single year (2023) is sufficient. Our methodology does not require a panel dataset
spanning multiple years'' \citep[p.~3]{lagasio2024} --- and its conclusion states the opposite: ``The
cross-sectional nature of our data limits causal inferences, highlighting the need for
longitudinal studies to examine how changes in corporate characteristics and
regulatory environments affect ESGSI over time'' \citep[p.~10]{lagasio2024}. The two sit in tension in
the original; our design resolves it in favour of the second. \S\ref{sec:temporal}
reports the resulting series, conditional on the measurement findings established
before it.

\paragraph{What the technical sections establish.} \S\ref{sec:original} audits what the
original specifies and what it leaves unspecified;
\S\ref{sec:data}--\S\ref{sec:index} document our corpus, extraction, lexicon and index.
\S\ref{sec:measurement} shows that \emph{vector} normalisation, not term weighting,
governs what a TF-IDF substance score measures, and \S\ref{sec:artefact} that the
original's headline conclusion turns on two undisclosed choices, one of them a
\emph{score} normalisation.
\S\ref{sec:validity} and \S\ref{sec:robustness} establish which substance
specification survives an independent validity check and how far the labels move
across specifications. \S\ref{sec:temporal} gives the 2018--2024 series,
\S\ref{sec:limitations} the limitations, and \S\ref{sec:repro} the full parameter,
vocabulary and lexicon listings.

\begin{table}[htbp]
\centering
\caption{Specification delta against \citet{lagasio2024}.}
\label{tab:scope-divergence}
\small
\begin{tabular}{@{}lp{4.6cm}p{5.2cm}@{}}
\toprule
 & \citet{lagasio2024} & This study \\
\midrule
Design & Cross-sectional, 2023 only & Balanced panel, 2018--2024 \\
Documents & 749 standalone ESG reports, one per firm & 343 unique reports: 235 annual reports, 101 URD, 7 standalone sustainability reports \\
Firms & 749, worldwide, from a Refinitiv top-2{,}000 starting sample & 49, seven European countries, each present in all seven years \\
Inclusion rule & English only; firms reporting ESG inside the annual report excluded & One annually recurring filing per firm-year, chiefly the annual report or URD --- largely the population the original excludes \\
Text vectorised & Whole report & Lexically extracted ESG zones (\S\ref{sec:extraction}) \\
ESG vocabulary & GRI/SASB-derived; size not stated, list not published & 154 terms, published in full (\S\ref{sec:repro}) \\
Tone dictionary & TextBlob polarity & Loughran--McDonald \\
Substance score & One score; aggregation over the vocabulary not specified & Three specifications, reported side by side (\S\ref{sec:index}) \\
Score normalisation & Prose says min--max; printed formula is a $z$-score difference & $z$-score; both examined (\S\ref{sec:artefact}) \\
\bottomrule
\end{tabular}
\end{table}
